% GENERATED FILE. Edit canonical lessons, then run npm run generate:books.
% canonical-source-hash: fnv1a64:e05ae37509f61c4f
% canonical-lessons: TE-C157-cellincu, TE-C157-veyincu, TE-C157-udikincu, TE-C157-kalupu, TE-C157-rubbu

\chapter{To pay, To fry, To boil, To mix, To grind}
\label{ch:te-to-pay-to-fry-to-boil-to-mix-to-grind}
\hlchaptermodality{157}

\begin{chapteropening}
\textbf{By the end of this chapter:} \emph{I can say to pay, to fry, to boil, to mix and to grind.}

Complete the last lesson of chapter 157: I can say to pay, to fry, to boil, to mix and to grind.
\end{chapteropening}

\section[\texorpdfstring{\te{చెల్లించు} (celliñcu)}{చెల్లించు (celliñcu)}]{\te{చెల్లించు} (celliñcu) — to pay}

\label{lesson:TE-C157-cellincu}



\begin{quote}
Before the new one: say the Telugu for an advertisement, then the Telugu for a receipt.
\end{quote}

\subsection*{You'll want to know: \te{చెల్లించు}}
\textbf{\te{చెల్లించు}} — \emph{celliñcu} — \textquotedblleft{}to pay\textquotedblright{}.

\begin{grammarlens}[title={What you've built}]
One more word for this chapter.
\end{grammarlens}

\subsection*{Your turn}
\begin{itemize}
\raggedright
  \item \emph{Say it:} \emph{celliñcu}
  \item \emph{Say it:} \emph{celliñcu}, once more
  \item \emph{Say it:} say \emph{celliñcu}
  \item \emph{Recall:} say the Telugu for an advertisement, then the Telugu for a receipt, then say \emph{celliñcu} again
\end{itemize}

\begin{tcolorbox}[breakable,colback=teal!4,colframe=teal!35!black,title={Before you move on}]
What does \te{చెల్లించు} mean? (\textquotedblleft{}To pay\textquotedblright{}.) Say it once more.
\end{tcolorbox}
\section[\texorpdfstring{\te{వేయించు} (vēyiñcu)}{వేయించు (vēyiñcu)}]{\te{వేయించు} (vēyiñcu) — to fry, to roast}

\label{lesson:TE-C157-veyincu}



\begin{quote}
Before the new one: say the Telugu for a receipt, then the Telugu for to pay.
\end{quote}

\subsection*{You'll want to know: \te{వేయించు}}
\textbf{\te{వేయించు}} — \emph{vēyiñcu} — \textquotedblleft{}to fry, to roast\textquotedblright{}.

\begin{grammarlens}[title={What you've built}]
One more word for this chapter.
\end{grammarlens}

\subsection*{Your turn}
\begin{itemize}
\raggedright
  \item \emph{Say it:} \emph{vēyiñcu}
  \item \emph{Say it:} \emph{vēyiñcu}, once more
  \item \emph{Say it:} say \emph{vēyiñcu}
  \item \emph{Recall:} say the Telugu for a receipt, then the Telugu for to pay, then say \emph{vēyiñcu} again
\end{itemize}

\begin{tcolorbox}[breakable,colback=teal!4,colframe=teal!35!black,title={Before you move on}]
What does \te{వేయించు} mean? (\textquotedblleft{}To fry, to roast\textquotedblright{}.) Say it once more.
\end{tcolorbox}
\section[\texorpdfstring{\te{ఉడికించు} (uḍikiñcu)}{ఉడికించు (uḍikiñcu)}]{\te{ఉడికించు} (uḍikiñcu) — to boil (food), to simmer}

\label{lesson:TE-C157-udikincu}



\begin{quote}
Before the new one: say the Telugu for to pay, then the Telugu for to fry.
\end{quote}

\subsection*{You'll want to know: \te{ఉడికించు}}
\textbf{\te{ఉడికించు}} — \emph{uḍikiñcu} — \textquotedblleft{}to boil (food), to simmer\textquotedblright{}.

\begin{grammarlens}[title={What you've built}]
One more word for this chapter.
\end{grammarlens}

\subsection*{Your turn}
\begin{itemize}
\raggedright
  \item \emph{Say it:} \emph{uḍikiñcu}
  \item \emph{Say it:} \emph{uḍikiñcu}, once more
  \item \emph{Say it:} say \emph{uḍikiñcu}
  \item \emph{Recall:} say the Telugu for to pay, then the Telugu for to fry, then say \emph{uḍikiñcu} again
\end{itemize}

\begin{tcolorbox}[breakable,colback=teal!4,colframe=teal!35!black,title={Before you move on}]
What does \te{ఉడికించు} mean? (\textquotedblleft{}To boil (food), to simmer\textquotedblright{}.) Say it once more.
\end{tcolorbox}
\section[\texorpdfstring{\te{కలుపు} (kalupu)}{కలుపు (kalupu)}]{\te{కలుపు} (kalupu) — to mix, to stir in}

\label{lesson:TE-C157-kalupu}



\begin{quote}
Before the new one: say the Telugu for to fry, then the Telugu for to boil.
\end{quote}

\subsection*{You'll want to know: \te{కలుపు}}
\textbf{\te{కలుపు}} — \emph{kalupu} — \textquotedblleft{}to mix, to stir in\textquotedblright{}.

\begin{grammarlens}[title={What you've built}]
One more word for this chapter.
\end{grammarlens}

\subsection*{Your turn}
\begin{itemize}
\raggedright
  \item \emph{Say it:} \emph{kalupu}
  \item \emph{Say it:} \emph{kalupu}, once more
  \item \emph{Say it:} say \emph{kalupu}
  \item \emph{Recall:} say the Telugu for to fry, then the Telugu for to boil, then say \emph{kalupu} again
\end{itemize}

\begin{tcolorbox}[breakable,colback=teal!4,colframe=teal!35!black,title={Before you move on}]
What does \te{కలుపు} mean? (\textquotedblleft{}To mix, to stir in\textquotedblright{}.) Say it once more.
\end{tcolorbox}
\section[\texorpdfstring{\te{రుబ్బు} (rubbu)}{రుబ్బు (rubbu)}]{\te{రుబ్బు} (rubbu) — to grind (wet batter)}

\label{lesson:TE-C157-rubbu}



\begin{quote}
Before the new one: say the Telugu for to boil, then the Telugu for to mix.
\end{quote}

\subsection*{You'll want to know: \te{రుబ్బు}}
\textbf{\te{రుబ్బు}} — \emph{rubbu} — \textquotedblleft{}to grind (wet batter)\textquotedblright{}.

\begin{grammarlens}[title={What you've built}]
One more word for this chapter.
\end{grammarlens}

\subsection*{Your turn}
\begin{itemize}
\raggedright
  \item \emph{Say it:} \emph{rubbu}
  \item \emph{Say it:} \emph{rubbu}, once more
  \item \emph{Say it:} say \emph{rubbu}
  \item \emph{Recall:} say the Telugu for to boil, then the Telugu for to mix, then say \emph{rubbu} again
\end{itemize}

\begin{tcolorbox}[breakable,colback=teal!4,colframe=teal!35!black,title={Before you move on}]
What does \te{రుబ్బు} mean? (\textquotedblleft{}To grind (wet batter)\textquotedblright{}.) Say it once more.
\end{tcolorbox}
